\documentclass[11pt]{article}
% Self-contained: compiles with a standard TeX distribution (no special class files).
\usepackage[utf8]{inputenc}
\usepackage[T1]{fontenc}
\usepackage{amsmath,amssymb}
\usepackage{booktabs}
\usepackage{multirow}
\usepackage{graphicx}
\usepackage{microtype}
\usepackage[margin=1in]{geometry}
\usepackage[hidelinks]{hyperref}
\usepackage{xcolor}
\usepackage{authblk}

\newcommand{\todo}[1]{\textcolor{red}{[#1]}}

\title{\bfseries Beyond Frame Accuracy: Online Surgical Phase Recognition\\
as Quickest Change Detection}
\author[1]{odafeng}
\affil[1]{\small (affiliation)}
\date{}

\begin{document}
\maketitle

\begin{abstract}
Frame-level accuracy on Cholec80 is saturated, yet it is a poor proxy for the
online deployability of a surgical phase recognition system: dominated by long,
stable phases, it is blind to how cleanly and promptly phase \emph{transitions}
are detected, and to whether the model's confidence can be trusted. We make three
contributions in the clinically realistic online/causal setting. \textbf{(i)} A
reproducible \emph{reliability evaluation suite}---relaxed accuracy, segmental
F1/edit, over-segmentation, transition-detection latency and false-alarms, and a
new \emph{boundary-region ECE} that exposes calibration failure localised at phase
transitions which standard ECE averages away. Applied identically across five
feature backbones and three causal temporal heads, it reveals failure modes that
are universal and invisible to accuracy: over-segmentation up to $10\times$ and
boundary-region calibration error $3$--$5\times$ the phase interior.
\textbf{(ii)} We \emph{reframe} online phase recognition as \emph{multi-phase
quickest change detection} (QCD). A CUSUM detector on the model's calibrated phase
posteriors, constrained by a phase-transition graph, enjoys a model-independent
structural bound on over-segmentation and---being asymptotically delay-optimal for
a given false-alarm rate---dominates hand-crafted post-processing on the
over-segmentation/latency frontier in $\sim$10/12 backbone-head cells, at a tunable
accuracy cost. \textbf{(iii)} A \emph{learned} causal sequential detector that
generalises the fixed rule reaches an operating point unreachable by fixed CUSUM or
heuristics---preserving frame accuracy \emph{and} low latency simultaneously. On a
different procedure (Cataract-101) the failure modes recur, confirming generality,
though their severity---and hence the benefit of QCD---is procedure dependent. We
report negative and nuanced results faithfully. Code and configurations are
released.
\end{abstract}

\section{Introduction}
Automated recognition of the surgical phase from laparoscopic video underpins
intra-operative decision support, automated documentation, and operating-room
workflow optimisation. On the de-facto benchmark Cholec80~\cite{twinanda2017},
offline relaxed-boundary accuracy now exceeds 95\%, and even our reproductions of
standard online heads reach $\sim$89\% video-averaged accuracy. Continued
score-chasing on a saturated metric yields little clinically actionable insight,
and the community has begun to question evaluation practice
itself~\cite{funke2023}.

We argue that the field should foreground \emph{reliability} in the
\emph{online/causal} setting (each frame is labelled using only the past)---the
only setting deployable during surgery. Two failure modes matter clinically and are
invisible to frame accuracy: (a) \textbf{transition behaviour}---a model correct on
90\% of frames may still thrash between phases around every transition, emitting
many spurious segments that frame accuracy, dominated by long phases, hides; and
(b) \textbf{trustworthy confidence}---for a model to defer or raise an alert, its
confidence must be calibrated \emph{where decisions are hard}, i.e. at transitions.

Our central observation is that these two quantities---detection delay and false
alarms at phase transitions---are precisely the operating characteristics of
\emph{quickest change detection} (QCD), a branch of sequential analysis with
decades of optimality theory~\cite{page1954,lorden1971,moustakides1986}. Reframing
online phase recognition as multi-phase change detection brings this theory to bear:
the messy ``post-process the predictions'' problem becomes one with provable
guarantees, and the clinically meaningful metrics become the objective itself.

\paragraph{Contributions.}
\begin{enumerate}\itemsep0pt
\item A reproducible online \emph{reliability suite} and the \emph{boundary-region
ECE} metric, applied identically across 5 backbones $\times$ 3 causal heads, that
exposes universal, accuracy-hidden failure modes (\S\ref{sec:suite},
\S\ref{sec:diagnosis}).
\item The \emph{QCD reframing} with a CUSUM/Shiryaev--Roberts decoder: a structural
over-segmentation bound, asymptotic delay-optimality, and an empirical domination of
heuristic post-processing, at a calibration-coupled, tunable accuracy cost
(\S\ref{sec:qcd}, \S\ref{sec:qcd-results}).
\item A \emph{learned sequential detector} that preserves accuracy and low latency
simultaneously (\S\ref{sec:learned}, \S\ref{sec:learned-results}); and a
cross-procedure study on Cataract-101 (\S\ref{sec:cross}).
\end{enumerate}

\section{Related Work}
\paragraph{Surgical phase recognition.} EndoNet/PhaseNet~\cite{twinanda2017}
established CNN per-frame models on Cholec80; recurrent and temporal-convolutional
models followed. TeCNO~\cite{czempiel2020} introduced a strictly causal multi-stage
temporal convolutional network for online use; transformer/Conformer heads such as
LoViT~\cite{liu2023} and aggregation transformers such as
Trans-SVNet~\cite{gao2021} model longer context. MS-TCN~\cite{farha2019} is the
offline multi-stage TCN we use as a reference. Most prior work reports frame
accuracy and per-phase precision/recall/Jaccard, frequently under a $\pm 10$\,s
relaxed boundary.

\paragraph{Evaluation critique.} Funke et al.~\cite{funke2023} showed that surgical
phase evaluation is inconsistent and that relaxed-boundary scoring inflates and
obscures differences. We extend that critique from accuracy-style metrics to
\emph{calibration} and \emph{online transition dynamics}, and pair it with a
principled decoder.

\paragraph{Calibration and quickest change detection.} Temperature
scaling~\cite{guo2017} is the standard post-hoc calibration; calibration has been
little studied for surgical phase recognition and, to our knowledge, never localised
to phase boundaries. Quickest change detection~\cite{page1954,lorden1971,moustakides1986}
studies the optimal trade-off between detection delay and
false-alarm rate; we instantiate it for the multi-phase, ordered surgical setting.

\section{A Reliability Evaluation Suite for Online Phase Recognition}
\label{sec:suite}
All metrics operate on per-video integer label streams at 1\,fps (so one frame
$=$ one second) and on per-frame probabilities; all are causal-safe.

\paragraph{Accuracy and segmental quality.} We report frame and video-averaged
accuracy, relaxed accuracy (within $\pm\tau$\,s of a ground-truth transition,
predicting either adjacent phase is not penalised), segmental
F1@$\{10,25,50\}$ and edit score~\cite{farha2019}, and an explicit
\emph{over-segmentation ratio} (predicted segments / ground-truth segments).

\paragraph{Transition latency.} For each ground-truth transition into phase $b$ at
time $t$, the detection latency is the first time $\ge t$ at which the prediction is
$b$ and remains $b$ for $\ge k$ frames; a switch to $b$ before $t$ is a false start.
Latency is computed only over \emph{clean} onsets (the model not already committed
to $b$), so early/flickered commitments do not masquerade as zero latency.

\paragraph{Boundary-region ECE.} Expected Calibration Error (ECE) averages over all
frames and is therefore \emph{diluted} by the many easy interior frames. We
partition frames into an \emph{interior} set and a \emph{boundary} set (within
$\pm\tau$\,s of any ground-truth transition) and report
\emph{interior-ECE} and \emph{boundary-region ECE} separately. The contrast
quantifies calibration failure localised exactly where the model is most likely to
err. We verify all conclusions are insensitive to $\tau\!\in\!\{5,10,15\}$\,s and
$k\!\in\!\{1,3,5\}$.

\section{Diagnosis: Accuracy Hides Failure}
\label{sec:diagnosis}
\paragraph{Setup.} We use the two-stage pipeline frame $\to$ frozen per-frame
feature $\to$ causal temporal head. Five backbones provide features: ImageNet
ResNet50; the surgical MAE model EndoViT~\cite{endovit} (pretrained, and fine-tuned
on Cholec80); an end-to-end fine-tuned ResNet50; and a self-trained masked
auto-encoder (\S\ref{sec:setup}). Three causal heads operate on identical features:
TeCNO~\cite{czempiel2020}, a causal LoViT~\cite{liu2023}, and a causal
ASFormer~\cite{yi2021}. Features are per-dimension standardised using training
statistics so that the temporal head is invariant to each backbone's feature scale.
Splits are the standard 32/8/40 train/val/test on Cholec80; the test set
(videos 41--80) is never used for training, tuning, calibration, or model selection.

\paragraph{Causality is airtight.} We added a prefix-invariance regression test: a
causal model's first-$k$ outputs must be identical whether the model is fed only the
first $k$ frames or the whole sequence. All three heads pass with
$\max|\Delta|\!\sim\!10^{-6}$. (This test caught a normalisation-over-time leak in an
early LoViT implementation, which we fixed.)

\paragraph{Finding.} Across all twelve backbone$\times$head baselines the
per-frame argmax stream over-segments severely and is badly mis-calibrated at
boundaries, while frame accuracy remains high. Online heads emit $\sim$5--$10\times$
more segments than ground truth, with segmental F1@10 around 25--40; boundary-region
ECE reaches 25--46\% while the interior is 5--9\% and global ECE looks benign
(6--11\%). Table~\ref{tab:diag} illustrates this for representative heads. The
phenomenon is a temporal-decoding problem, not a feature-quality one: it persists at
every backbone strength.

\begin{table}[t]\centering\small
\caption{Accuracy hides failure (Cholec80 test, standardised ResNet50 features,
illustrative). Frame accuracy is high while over-segmentation and boundary-region
ECE are severe.}
\label{tab:diag}
\begin{tabular}{lccccc}
\toprule
head (raw argmax) & acc \% & over-seg & F1@10 & interior-ECE & boundary-ECE \\
\midrule
TeCNO    & 88.6 & 6.7$\times$ & 32.5 & 7.2 & 31.7 \\
LoViT    & 87.5 & 9.5$\times$ & 25.8 & 9.3 & 35.5 \\
ASFormer & 87.9 & 6.4$\times$ & 37.3 & 6.0 & 33.6 \\
\bottomrule
\end{tabular}
\end{table}

\section{Online Phase Recognition as Quickest Change Detection}
\label{sec:qcd}
\paragraph{Setup.} Let $p_t=P(z_t\mid x_{1:t})$ be the causal posterior over $K$
phases produced by the temporal head. The true phase $z_t$ is piecewise constant
with transitions; Cholec80 phases follow a partial order with limited revisits,
encoded as a transition graph $G$ estimated from the \emph{training} labels only.

\paragraph{Detector.} Maintain a declared phase $\hat k$. For each allowed successor
$j$ of $\hat k$ in $G$, accumulate the posterior log-likelihood ratio with a CUSUM
recursion
\begin{equation}
W^{j}_t=\max\!\big(0,\; W^{j}_{t-1}+\log p_t[j]-\log p_t[\hat k]\big),\qquad
\text{declare } \hat k\!\to\! j^\star \text{ when } \max_j W^{j}_t \ge h ,
\end{equation}
then reset. The threshold $h$ trades detection delay against false-alarm rate; a
Shiryaev--Roberts variant is also implemented.

\paragraph{Two guarantees.}
\emph{(1) Structural over-segmentation bound (model-independent).} The number of
predicted segments equals the number of declared transitions; with the ordered
graph $G$ (a DAG plus $R$ allowed revisits), forward progress bounds the segment
count to $\le (K{-}1)+R$, so the over-segmentation ratio is $\approx 1$ by
construction---a guarantee no frame-wise classifier or EMA smoother provides. On
synthetic streams a per-frame argmax emits 189 segments where CUSUM emits 3.
\emph{(2) Delay optimality.} Under correct specification the per-transition CUSUM is
asymptotically minimax in Lorden's sense---minimising worst-case expected detection
delay subject to a false-alarm constraint~\cite{lorden1971,moustakides1986}. This is
exactly the frontier on which hand-crafted smoothing is sub-optimal.

\paragraph{Calibration coupling.} The log-likelihood ratio is meaningful only if
posteriors are calibrated; mis-calibration biases the increment and degrades
delay/false-alarm. This makes temperature scaling / deep-ensemble calibration part
of the method, and makes boundary-region ECE the right diagnostic for where the
detector is most fragile.

\section{A Learned Sequential Detector}
\label{sec:learned}
The fixed-rule CUSUM is optimal only under correct specification and pays an
accuracy cost through the hard ordering constraint. We therefore \emph{learn} a
causal detector that generalises the rule. Maintain a soft state
$s_t=(1-g_t)s_{t-1}+g_t\,p_t$ with a learned gain $g_t\in(0,1)$ produced by a small
causal GRU from $[p_t,s_{t-1},\mathrm{conf}_t,\mathrm{entropy}_t]$ and an evidence-
accumulating hidden state. It is trained on a differentiable surrogate of the
delay/false-alarm trade-off,
\begin{equation}
\mathcal{L}=\underbrace{\mathrm{CE}(s_t,z_t)}_{\text{tracking (delay/accuracy)}}
+\;\lambda\underbrace{\textstyle\frac1T\sum_t\lVert s_t-s_{t-1}\rVert_1}_{\text{smoothness (false alarm / over-seg)}},
\end{equation}
and sweeping $\lambda$ traces the operating curve. The detector is strictly causal
and streaming-deployable; it is trained on training-video posteriors, selected on
val, and evaluated on test (a learned second-stage decoder; no leakage). The network
learns to open the gate at genuine transitions and close it in phase interiors,
which fixed rules cannot do adaptively.

\section{Experiments}
\subsection{Setup}
\label{sec:setup}
Cholec80, 32/8/40 train/val/test. Five backbones produce frozen, standardised
features; three causal heads (TeCNO, LoViT, ASFormer) are trained with five seeds
each (the seeds double as a deep ensemble). The self-trained MAE backbone is a
ViT-B/16 pretrained from scratch with masked auto-encoding on training+validation
frames only (videos 1--40; the test set is never seen, removing the contamination
caveat that affects publicly pretrained surgical backbones). Cross-procedure
evaluation uses Cataract-101~\cite{cataract101} (101 cataract surgeries, 10 phases).

\subsection{Baselines reproduce prior accuracy}
Our online TeCNO reaches 88.96\% video-averaged accuracy (prior 88.95\%) and offline
MS-TCN 90.84\% (prior 90.81\%), confirming a faithful pipeline. Table~\ref{tab:bb}
reports the standardised five-backbone baseline accuracies.

\begin{table}[t]\centering\small
\caption{Baseline frame accuracy (\%) on Cholec80 test (standardised features,
5-seed mean). End-to-end and fine-tuned backbones are strongest; the self-trained
single-dataset MAE is the weakest---an honest, contamination-free control.}
\label{tab:bb}
\begin{tabular}{lccc}
\toprule
backbone & TeCNO & LoViT & ASFormer \\
\midrule
ResNet50 (ImageNet)        & 88.6 & 87.5 & 87.9 \\
EndoViT (pretrained)       & 82.2 & 79.1 & 76.4 \\
EndoViT (fine-tuned)       & 88.7 & 87.1 & 87.3 \\
End-to-end ResNet50        & 88.9 & 90.1 & 89.1 \\
Self-trained MAE (ours)    & 59.0 & 28.2 & 39.4 \\
\bottomrule
\end{tabular}
\end{table}

\subsection{QCD dominates the segmentation/latency frontier}
\label{sec:qcd-results}
At matched over-segmentation ($\le 2\times$), the CUSUM decoder achieves lower
detection latency than the best heuristic post-processor in $\sim$10/12
backbone$\times$head cells (ties only on the weak EndoViT-pretrained backbone), and
reaches the structural over-segmentation floor ($\sim$1.2--1.8$\times$).
Table~\ref{tab:qcd} gives representative cells. The ordering graph provides a
\emph{tunable accuracy$\leftrightarrow$segmentation knob}: with the estimated graph,
accuracy drops (the hard constraint overrides the model when its evidence
contradicts the learned order), but a relaxed/no-ordering graph recovers accuracy to
$\sim$87--90\% at the cost of cleaner segmentation. We report this trade-off
explicitly rather than hiding it.

\begin{table}[t]\centering\small
\caption{QCD vs.\ heuristic post-processing at over-seg $\le 2$ (Cholec80 test).
CUSUM reaches lower latency but pays an accuracy cost via the ordering graph; the
heuristic keeps accuracy at higher latency.}
\label{tab:qcd}
\begin{tabular}{llccc}
\toprule
cell & decoder & over-seg & latency (s) & acc \% \\
\midrule
\multirow{2}{*}{ResNet50/TeCNO} & heuristic & 1.8 & 10.7 & 88.5 \\
                                & CUSUM     & 1.8 & \textbf{7.9} & 79.3 \\
\multirow{2}{*}{ResNet50/LoViT} & heuristic & 1.9 & 20.8 & 87.7 \\
                                & CUSUM     & 1.9 & \textbf{10.9} & 71.1 \\
\multirow{2}{*}{e2e/LoViT}      & heuristic & 1.8 & 15.6 & 89.6 \\
                                & CUSUM     & 1.9 & \textbf{7.3} & 74.7 \\
\bottomrule
\end{tabular}
\end{table}

\paragraph{Calibration coupling.} Better-calibrated posteriors improve QCD, most
visibly on the weak LoViT head: a 5-seed deep ensemble raises QCD frame accuracy
from 71.1\% to 76.8\% and F1@10 from 63 to 72, while a single global temperature
already reduces ECE. This confirms the calibration--detection coupling predicted by
the theory.

\subsection{The learned detector preserves accuracy \emph{and} low latency}
\label{sec:learned-results}
No single fixed decoder achieves all of \{low over-seg, low latency, high accuracy\}:
CUSUM minimises over-seg and latency but sacrifices accuracy (the ordering tax);
heuristics keep accuracy but pay latency. The learned detector is the only decoder
that simultaneously preserves frame accuracy \emph{and} low latency while
substantially cutting over-segmentation (Table~\ref{tab:learned}). It does not reach
the absolute minimum over-segmentation of CUSUM, but occupies the
clinically-preferred high-accuracy/low-latency corner unreachable by fixed rules.

\begin{table}[t]\centering\small
\caption{Decoders on e2e/LoViT (Cholec80 test). The learned detector keeps argmax-
level accuracy at low latency; CUSUM sacrifices accuracy, the heuristic pays latency.}
\label{tab:learned}
\begin{tabular}{lcccc}
\toprule
decoder & over-seg & acc \% & latency (s) & F1@10 \\
\midrule
argmax (raw)         & 4.4 & 90.7 & 6.3 & 42.9 \\
heuristic ($\le$2)   & 1.8 & 89.8 & 13.8 & 74.1 \\
CUSUM ($\le$2)       & 1.6 & 74.0 & 6.9 & 70.1 \\
\textbf{learned} ($\lambda{=}64$) & 2.5 & \textbf{90.4} & \textbf{9.1} & 63.2 \\
\bottomrule
\end{tabular}
\end{table}

\subsection{Cross-procedure generalisation (Cataract-101)}
\label{sec:cross}
We replicate the full pipeline on Cataract-101, a procedure utterly different from
cholecystectomy (eye microsurgery, 10 phases). A TeCNO trained on cataract reaches
92.3\% validation accuracy. The same failure modes recur---over-segmentation
$2.9\times$ and boundary-region ECE 14.6\%, hidden by 87.4\% frame
accuracy---confirming that the phenomenon is general and not a Cholec80 artefact.
However, both are \emph{milder} than on Cholec80: cataract has short phases and
sharp transitions, so detection latencies are already tiny (1--3\,s vs.\ 6--20\,s)
and over-segmentation already low. Consequently the QCD benefit is smaller here,
while the CUSUM accuracy cost persists (87.4\%$\to$83.7\%). The honest reading is
that the failure modes generalise but their \emph{severity}---and hence the payoff
of principled decoding---is procedure dependent, largest on messy, long-phase
procedures.

\section{Discussion and Limitations}
We report results faithfully, including negative and nuanced ones.
\textbf{(i) Method strength.} The fixed CUSUM dominates heuristics on the
segmentation/latency frontier but at an accuracy cost; the learned detector removes
that cost on the accuracy/latency axes but does not reach the absolute minimum
over-segmentation. There is no decoder that wins on all three axes simultaneously;
the contribution is making the trade-off explicit and learnable.
\textbf{(ii) Statistics.} Accuracy ``non-inferiority'' should be established with a
pre-registered margin and a TOST/confidence-interval procedure rather than a
non-significant two-sided test; we report paired deltas and intend bootstrap
confidence intervals and family-wise correction across the metric grid in the
camera-ready.
\textbf{(iii) Backbones.} Our self-trained single-dataset MAE yields weak frozen
features (a known property of MAE, which favours fine-tuning); it serves as an
honest contamination-free control rather than a state-of-the-art backbone.
\textbf{(iv) Datasets.} Cholec80 is the primary dataset; Cataract-101 provides one
cross-procedure point. We attempted to obtain AutoLaparo (hysterectomy) but its
distribution link has been unreachable for several months.
\textbf{(v) Heads.} We compare three causal heads; integrating recent baselines
(Trans-SVNet, Surgformer) is ongoing and only adds posterior sources to the same
analysis.

\section{Conclusion}
High frame accuracy is not reliability. Measuring online phase recognition with a
boundary-localised, transition-aware suite exposes universal failure modes that
accuracy hides; reframing the problem as quickest change detection brings sequential
optimality and a structural over-segmentation guarantee; and a learned sequential
detector reaches a high-accuracy, low-latency operating point that fixed rules
cannot. The benefits are largest precisely where surgery is messiest.

\begin{thebibliography}{99}\small
\bibitem{twinanda2017} A.~P. Twinanda et al. EndoNet: A deep architecture for
recognition tasks on laparoscopic videos. \emph{IEEE TMI}, 2017.
\bibitem{czempiel2020} T.~Czempiel et al. TeCNO: Surgical phase recognition with
multi-stage temporal convolutional networks. \emph{MICCAI}, 2020.
\bibitem{farha2019} Y.~A. Farha and J.~Gall. MS-TCN: Multi-stage temporal
convolutional network for action segmentation. \emph{CVPR}, 2019.
\bibitem{liu2023} Y.~Liu et al. LoViT: Long video transformer for surgical phase
recognition. \emph{Medical Image Analysis}, 2023.
\bibitem{gao2021} X.~Gao et al. Trans-SVNet: Accurate phase recognition from
surgical videos via hybrid embedding aggregation transformer. \emph{MICCAI}, 2021.
\bibitem{yi2021} F.~Yi et al. ASFormer: Transformer for action segmentation.
\emph{BMVC}, 2021.
\bibitem{funke2023} I.~Funke et al. Metrics Matter in Surgical Phase Recognition.
\emph{MICCAI}, 2023.
\bibitem{guo2017} C.~Guo et al. On calibration of modern neural networks.
\emph{ICML}, 2017.
\bibitem{page1954} E.~S. Page. Continuous inspection schemes. \emph{Biometrika},
1954.
\bibitem{lorden1971} G.~Lorden. Procedures for reacting to a change in distribution.
\emph{Ann. Math. Statist.}, 1971.
\bibitem{moustakides1986} G.~V. Moustakides. Optimal stopping times for detecting
changes in distributions. \emph{Ann. Statist.}, 1986.
\bibitem{endovit} S.~Batic et al. (EndoViT) Pretraining vision transformers on
surgical data. (cf. egeozsoy/EndoViT).
\bibitem{cataract101} K.~Schoeffmann et al. Cataract-101: video dataset of 101
cataract surgeries. \emph{ACM MMSys}, 2018.
\end{thebibliography}

\end{document}
